\documentclass[11pt]{article}
\usepackage[utf8]{inputenc}
\usepackage[T1]{fontenc}
\usepackage[a4paper,margin=3.1cm]{geometry}
\usepackage{newtxtext,newtxmath}
\usepackage[hidelinks]{hyperref}
\pagestyle{empty}
\begin{document}
\begin{center}
{\fontsize{17}{20}\selectfont The Value of the Marginal Solver:\\[2pt]
Competition and Best Execution in Intent-Based Markets\par}
\vspace{0.8em}
{\fontsize{11}{14}\selectfont Nils Bundi\par}
{\fontsize{9}{11}\selectfont Zurich University of Applied Sciences (ZHAW), \texttt{bund@zhaw.ch}\par}
\end{center}
\vspace{1.2em}
\noindent\textbf{Abstract.}\quad Intent-based exchanges, the fastest-growing execution model in on-chain finance, replace order routing with a public solver auction. CoW Protocol is, to my knowledge, the first major exchange to publish the full competition behind each trade, winning and losing bids alike. I use it to measure the value of one more competitor for a customer order. From its 565{,}259 auctions (twelve months, 692{,}964 orders, 47 solvers), I find that a single retained bidder costs the median order 14 basis points, a second restores 8, a fifth under 1. The marginal-solver value curve shifts upward as liquidity thins. A reward reform during the sample period offers a natural experiment: three major solvers exited the field, and rivals filled their place within a month with no impact on bids. Solver competition is supplied by the field's depth, not its members' identity. I conclude with the observation that where a broker is bound to best execution by regulation, intent exchanges deliver it through competition that anyone can measure.

\vspace{1.4em}
\noindent\textbf{Keywords:}\quad Solver competition; order flow auctions; best execution; batch auctions; market microstructure.
\end{document}
